\section{Framework}
\label{sec:framework}

The framework transforms existing scientific knowledge into ranked,
falsifiable questions through six stages
(Figure~\ref{fig:architecture}). Throughout, the design principle is
\emph{traceability}: every artifact carries the provenance needed to
audit or reproduce it.

\begin{figure}[t]
  \centering
  \todo{architecture figure: Existing Scientific Knowledge $\to$
  Evidence Representation $\to$ Scientific Reasoning Engine
  (Contradictions / Gaps / Alternative Explanations) $\to$
  Evidence-based Question Refinement $\to$ Ranked High-value Questions}
  \caption{Architecture overview. Signals flow from a scope-controlled
  corpus through claim extraction and a typed evidence graph into
  question generation, falsification screening, and two-stage ranking.}
  \label{fig:architecture}
\end{figure}

\subsection{Evidence representation}
\label{sec:framework:evidence}

Knowledge enters as three layers, collected in order of increasing cost:
a \emph{literature} layer (what has been claimed), a \emph{catalog}
layer (what objects exist, with parameters and uncertainties), and an
\emph{observation} layer (where the underlying data lives---metadata
only; data products are fetched on demand after a question is selected).
Every raw record is stored immutably inside a provenance envelope
$(\textit{source}, \textit{retrieved-at}, \textit{query-version},
\textit{source-id}, \textit{payload})$, then normalized into relational
tables.

The atomic unit is the \textbf{claim}: a tuple
\[
c = (\textit{text},\ \textit{evidence},\ \textit{assumptions},\
     \textit{uncertainty},\ \textit{objects},\ \textit{datasets},\
     \textit{location},\ \textit{tier})
\]
where \textit{location} points back into the source document and
\textit{tier} records the verification provenance:
\texttt{human\_gold} (annotated by a domain reviewer),
\texttt{model\_matched\_gold} (model extraction semantically matching a
gold claim; the canonical text remains the human's, both provenances are
kept), or \texttt{model\_only} (unreviewed model extraction, default
confidence 0.75). Model claims that duplicate gold claims are absorbed
into a single canonical node rather than added as siblings---duplicate
nodes would both inflate the graph and manufacture spurious tension.

\subsection{Evidence graph and typed tensions}
\label{sec:framework:graph}

Claims, objects, datasets, and assumptions become nodes; deterministic
rules add edges (\textit{uses-dataset}, \textit{measures},
\textit{assumes}), each recording the rule that created it. Cross-paper
claim pairs that share an object \emph{and} a topic term become
\emph{tension candidates}---deliberately a weak, high-recall rule. The
key representational commitment is that candidate tensions are
\textbf{typed by human adjudication} before they may drive question
generation, using the vocabulary
\[
\{\textit{contradicts},\ \textit{qualifies},\
  \textit{challenges-method},\ \textit{explains-discrepancy},\
  \textit{complements},\ \textit{supports}\}.
\]
This vocabulary is an empirical result of the calibration review
(Section~\ref{sec:experiments:tension}): most high-scoring candidates
are not contradictions, and a graph that collapses ``adds limiting
conditions'' into ``negates'' will amplify false conflicts as it grows.

\subsection{Question generation}
\label{sec:framework:questions}

Questions are generated only from prioritized signal classes:
(P1) human-confirmed observational tensions, (P2) methodological
challenges, (P3) qualifications across independent datasets, and
(P4) single-dataset conclusions from trusted claims. Unreviewed tension
candidates are excluded by design. Signal assembly is deterministic;
a language model only phrases the question, and every candidate carries
its evidence trail (claim identifiers, papers, adjudication basis) plus
an archival-data availability count---the first falsification screen.
Human editorial decisions (merges into question families, rewrites,
rejections) are recorded as committed curation records and re-applied
deterministically.

\subsection{Two-stage ranking}
\label{sec:framework:ranking}

Stage~A is a quality gate: evidence grounding plus a six-criterion
clarity rubric (presupposed conclusions, statistics/physics conflation,
explicit comparison, falsifiable outcome, bounded scope, stateable
failure condition). Clarity gates rather than scores---in calibration it
scored uniformly and only added a constant. Human-curated questions
bypass the model gate, whose verdict is retained as advisory.

Stage~B ranks survivors by
\[
\textit{scientific priority} =
0.35\,S + 0.25\,T + 0.20\,F + 0.10\,N + 0.10\,G,
\]
where $S$ (significance) is judged by a model but hard-capped by signal
tier (P1:~10, P2:~9, P3:~8.5, P4:~7.5) so that single-dataset robustness
checks cannot rank beside confirmed tensions; $T$ (tension strength) is
structural; $F$ (feasibility) averages judged subcomponents
(availability of the \emph{right} data rather than raw record counts,
data independence, analysis readiness, reanalysis sufficiency);
$N$ (novelty) is deliberately down-weighted---the system generates
questions from existing evidence tensions, so surface novelty is not the
point---and blends a judge with a corpus-similarity check; $G$ is
expected information gain. \emph{Execution priority} (feasibility) is
reported separately: abundant archives must not outrank a confirmed
cross-instrument tension, and a question can be scientifically first
while operationally harder.
